\documentclass[12pt,a4paper]{article}
\usepackage[utf8]{inputenc}
\usepackage[T1]{fontenc}
\usepackage{babel}
\usepackage{hyperref}
\usepackage{listings}
\usepackage{xcolor}

\definecolor{codegreen}{rgb}{0,0.6,0}
\definecolor{codegray}{rgb}{0.5,0.5,0.5}
\definecolor{codepurple}{rgb}{0.58,0,0.82}
\definecolor{backcolour}{rgb}{0.95,0.95,0.92}
\definecolor{keyword}{rgb}{0,0,1}
\definecolor{string}{rgb}{0.58,0,0.82}
\definecolor{comment}{rgb}{0.57,0.57,0.57}
\definecolor{type}{rgb}{0.1,0.5,0.1}

\lstdefinestyle{cstyle}{
    backgroundcolor=\color{backcolour},
    commentstyle=\color{comment}\textit,
    keywordstyle=\color{keyword}\bfseries,
    stringstyle=\color{string},
    numberstyle=\tiny\color{codegray},
    basicstyle=\ttfamily\small,
    breaklines=true,
    frame=single,
    framesep=5pt,
    frameround=tttt,
    numbers=left,
    numbersep=8pt,
    captionpos=b,
    tabsize=4,
    showstringspaces=false
}

\lstset{style=cstyle}

\title{Structures de Configuration GTK4 \\ Rapport Technique}
\author{Projet GTK Wrapper}
\date{\today}

\begin{document}

\maketitle

\tableofcontents

\newpage

\section{Introduction}
Ce rapport documente les structures de configuration de l'API GTK Wrapper. L'approche utilisée repose sur la configuration déclarative : au lieu d'appeler plusieurs fonctions, le développeur remplit une structure de configuration puis la passe à une fonction de création.

Ce document détaille les structures pour les widgets du wrapper avec une description de chaque champ et des exemples d'utilisation.

\newpage

\section{Structure \texttt{window\_config}}

La structure \texttt{window\_config} permet de configurer une fenêtre GTK4. Elle regroupe les paramètres géométriques, visuels et comportementaux nécessaires à sa création.

\subsection{Définition}

\begin{lstlisting}[language=C]
typedef struct {
    const char *title;
    const char *icon_name;
    const char *widget_name;
    const char *background_color;
    const char *background_image_path;
    int default_width;
    int default_height;
    int min_width;
    int min_height;
    int max_width;
    int max_height;
    bool resizable;
    bool decorated;
    bool modal;
    bool maximized;
    bool present_on_create;
    widget_style_config style;
} window_config;
\end{lstlisting}

\subsection{Description des champs}

\begin{center}
\begin{tabular}{|l|l|p{5cm}|}
\hline
\textbf{Champ} & \textbf{Type} & \textbf{Description} \\
\hline
title & const char* & Titre affiché dans la barre de titre \\
\hline
icon\_name & const char* & Icône du thème pour la barre de titre \\
\hline
widget\_name & const char* & Nom interne pour le débogage \\
\hline
background\_color & const char* & Couleur de fond (format CSS) \\
\hline
background\_image\_path & const char* & Chemin vers une image de fond \\
\hline
default\_width & int & Largeur par défaut en pixels \\
\hline
default\_height & int & Hauteur par défaut en pixels \\
\hline
min\_width & int & Largeur minimale \\
\hline
min\_height & int & Hauteur minimale \\
\hline
max\_width & int & Largeur maximale \\
\hline
max\_height & int & Hauteur maximale \\
\hline
resizable & bool & Autorise le redimensionnement \\
\hline
decorated & bool & Affiche la barre de titre \\
\hline
modal & bool & Fenêtre modale \\
\hline
maximized & bool & État maximisé à la création \\
\hline
present\_on\_create & bool & Affichage immédiat \\
\hline
style & widget\_style\_config & Configuration du style \\
\hline
\end{tabular}
\end{center}

\bigskip

Le champ \textbf{title} définit le texte affiché dans la barre de titre. Utilisé par les lecteurs d'écran pour l'accessibilité. La valeur peut être NULL.

Le champ \textbf{icon\_name} spécifie l'icône à afficher. Accepte les noms du thème système ("dialog-information", "document-open"). NULL utilise l'icône par défaut.

Le champ \textbf{widget\_name} définit un nom interne pour le débogage et les styles CSS. Apparaît dans les outils de développement GTK.

Le champ \textbf{background\_color} spécifie une couleur de fond au format CSS : noms ("red", "blue"), hexadécimal ("\#FF0000") ou RGB ("rgb(255, 0, 0)"). NULL utilise le thème par défaut.

Le champ \textbf{background\_image\_path} definit une image de fond. Le chemin doit pointer vers un fichier image valide. NULL desactive cette fonctionnalite.

Les champs \textbf{default\_width} et \textbf{default\_height} définissent les dimensions initiales en pixels.

Les champs \textbf{min\_width}, \textbf{min\_height}, \textbf{max\_width} et \textbf{max\_height} définissent les contraintes de dimensionnement. 0 désactive la limite.

Le champ \textbf{resizable} contrôle si l'utilisateur peut redimensionner la fenêtre. FALSE désactive les poignées de redimensionnement.

Le champ \textbf{decorated} contrôle la présence de la barre de titre et des bordures. FALSE crée une fenêtre sans barre de titre.

Le champ \textbf{modal} crée une fenêtre modale bloquant l'interaction avec les autres fenêtres jusqu'à sa fermeture.

Le champ \textbf{maximized} spécifie si la fenêtre est maximisée dès sa création.

Le champ \textbf{present\_on\_create} contrôle l'affichage immédiat après création.

Le champ \textbf{style} est une structure \texttt{widget\_style\_config} pour configurer le style.

\subsection{Exemple d'utilisation}

\begin{lstlisting}[language=C]
#include "window.h"

void create_main_window(GtkApplication *app) {
    window_config config = {
        .title = "Mon Application GTK",
        .default_width = 1024,
        .default_height = 768,
        .min_width = 800,
        .min_height = 600,
        .resizable = true,
        .decorated = true,
        .modal = false,
        .maximized = false,
        .present_on_create = true,
        .style = {
            .css_class = "main-window",
            .hexpand = true,
            .vexpand = true
        }
    };

    GtkWidget *window = create_window(app, &config);
    gtk_application_add_window(app, GTK_WINDOW(window));
}
\end{lstlisting}

\section{Structure \texttt{button\_config}}

La structure \texttt{button\_config} permet de configurer les boutons GTK4. Elle définit le label, l'icône, le style et le comportement callback du bouton.

\subsection{Définition}

\begin{lstlisting}[language=C]
typedef struct {
    const char *label;
    const char *icon_name;
    const char *theme_variant;
    bool use_underline;
    bool has_frame;
    bool can_shrink;
    int icon_size;
    widget_style_config style;
    void (*on_clicked)(GtkButton *button, gpointer user_data);
    gpointer user_data;
} button_config;

typedef struct {
    const char *label;
    GtkWidget *group_with;
    bool is_active;
    widget_style_config style;
    void (*on_toggled)(GtkCheckButton *button, gpointer user_data);
    gpointer user_data;
} radio_button_config;
\end{lstlisting}

\subsection{Description des champs}

\begin{center}
\begin{tabular}{|l|l|p{5cm}|}
\hline
\textbf{Champ} & \textbf{Type} & \textbf{Description} \\
\hline
label & const char* & Texte affiché sur le bouton \\
\hline
icon\_name & const char* & Icône du thème \\
\hline
theme\_variant & const char* & Variante de thème CSS \\
\hline
use\_underline & bool & Active les mnémoniques \\
\hline
has\_frame & bool & Affiche la bordure \\
\hline
can\_shrink & bool & Autorise la réduction \\
\hline
icon\_size & int & Taille de l'icône \\
\hline
style & widget\_style\_config & Configuration du style \\
\hline
on\_clicked & callback & Fonction appelée au clic \\
\hline
user\_data & gpointer & Données utilisateur \\
\hline
\end{tabular}
\end{center}

\bigskip

Le champ \textbf{label} définit le texte affiché sur le bouton. Peut contenir des mnémoniques ("\_Fichier" pour Alt+F). NULL pour un bouton purement iconique.

Le champ \textbf{icon\_name} spécifie une icône du thème système ("document-open", "edit-copy", etc.). Peut être combiné avec le label.

Le champ \textbf{theme\_variant} applique une variante : "primary" (principal), "success" (positif), "danger" (destructif), "secondary" (secondaire).

Le champ \textbf{use\_underline} active le traitement des mnémoniques clavier.

Le champ \textbf{has\_frame} contrôle l'affichage de la bordure. FALSE crée un bouton plat.

Le champ \textbf{can\_shrink} permet au bouton de réduire sa taille minimale.

Le champ \textbf{icon\_size} définit la taille : GTK\_ICON\_SIZE\_BUTTON (24x24), GTK\_ICON\_SIZE\_SMALL\_TOOLBAR (16x16), etc.

Le champ \textbf{style} est une structure \texttt{widget\_style\_config} pour le style.

Le champ \textbf{on\_clicked} est le callback appelé lors du clic.

Le champ \textbf{user\_data} transmet des données au callback.

\subsection{Boutons radio}

La structure \texttt{radio\_button\_config} permet de créer des boutons radio pour les selections exclusives.

Le champ \textbf{group\_with} lie le bouton à un groupe existant.

Le champ \textbf{is\_active} définit l'état sélectionné initial.

Le champ \textbf{on\_toggled} est appelé lors du changement d'état.

\subsection{Exemple d'utilisation}

\begin{lstlisting}[language=C]
#include "button.h"
#include <stdio.h>

void on_save_clicked(GtkButton *button, gpointer user_data) {
    g_print("Bouton clique!\n");
}

void create_ui_example(GtkWidget *box) {
    button_config save_btn = {
        .label = "_Sauvegarder",
        .theme_variant = "primary",
        .has_frame = true,
        .on_clicked = on_save_clicked,
        .user_data = NULL
    };
    GtkWidget *save_button = create_button(&save_btn);
    gtk_box_append(GTK_BOX(box), save_button);

    radio_button_config opt1 = {
        .label = "Option A",
        .is_active = true
    };
    GtkWidget *radio1 = create_radio_button(&opt1);
    gtk_box_append(GTK_BOX(box), radio1);

    radio_button_config opt2 = {
        .label = "Option B",
        .group_with = radio1
    };
    GtkWidget *radio2 = create_radio_button(&opt2);
    gtk_box_append(GTK_BOX(box), radio2);
}
\end{lstlisting}


\end{document}
